\scanpage{306}
\mspage{5 \deleted{\unsure{5}}}
Ce n'est autre que le foncteur représenté par $K_*$, et on le notera aussi
\[
 \deleted{\unsure{K_*(X_*)\longmapsto|K_*(X_*)=}}\qquad
 X_*\longmapsto K_*(X_*)=\mathrm{Hom}(X_*,K_*)
\]
Si $B=(\mathrm{Ab})$, $K_*$ sera de $=$ un objet de $\mathrm{Ss}_{\mathrm{ab}}$ \deleted{\unsure{$\mathrm{Ab}_*$}}, et en effet le foncteur qu'il représente est automatiquement muni de \deleted{\unsure{telle str.}} \inserted{\unsure{'une structure|type structure|une structure}} \unsure{abélien|abélienne}, $\overline{M}$ remarque pour tte structure algébrique (sur un ou plusieurs objets de base) qui ``puisse s'exprimer en termes de $\varprojlim$ exclusivement'' (grps, \deleted{\unsure{unités|modules}} anneaux, modules sur \deleted{\unsure{des algèbres, etc|tels, etc}} etc).

\uline{Rem} Les \deleted{foncteurs} invariants ainsi obtenus à coups de foncteurs représentables -- plus gén, d'objets simpliciaux ou cosimpliciaux de catégories $B$ ou $A$ -- sont de nature trop ``grosse'' et ``fruste'' pour être intéressants directement -- \deleted{\unsure{en}} en particulier, ce ne sont pas des invariants du type d'homotopie de $X_*$. On devra en extraire des invariants plus subtils qui soient des invariants du type d'homotopie de $X_*$, -- ceux-ci ne seront pas de nature si simple -- i.e. ``repr.'' par des objets simpliciaux ou cosimpliciaux. Notre propos sera de mettre sur les invariants ``gros'' qu'on vient de décrire, un maximum de structures supplémentaires \unsure{raisonnables}, et dans les catégories de structures de ce type (complexes de chaines ou de cochaines, algèbres cosimpliciales etc) d'inverser les flèches qu'on obtient à partir d'équiv. d'homotopie des $\mathrm{Ss}$, et de passer à des ``catégories de fractions'' en inversant ces flèches. L'invariant ainsi obtenu \deleted{\unsure{donne certes}} \inserted{(plus ``grossier''} bien sûr, mais plus subtil dans sa définition, et moins ``redondant'') exprimera alors de façon plus ou moins complète le type d'homotopie, ou en saisira de façon adéquate tels ou tels aspects.

Dans le cas où $A$ est une catégorie additive karoubienne (tout projecteur sur un objet a un noyau, ou ce qui revient au $=$, tt proj. a un conoyau) le \unsure{topoïs|topos|topo|théorie} de Dold-Puppe donne des équivalences canoniques
\diagram{\begin{tikzcd}[ampersand replacement=\&]
 A_* \arrow[r,shift left=1.1ex,"{\substack{\mathrm{ND}_{.}\\\approx}}"]
 \& A_{.}\arrow[l,shift left=1.1ex,"{\approx}"',"{\mathrm{DP}_*}"] \& {,}
\end{tikzcd}}
\margin{$\mathrm{ND}=$ ``non dégénéré''\\
$\mathrm{DP}=$ ``Dold-Puppe''}
où $A_.$ désigne la catégorie des \uline{complexes de chaines} dans $A$. On désigne par une $=$ lettre \inserted{telle $K$} deux objets qui se correspondent, mais en mettant en indice soit $*$ soit $.$, suivant qu'il
% archive sheet number: 307
